% !TEX TS-program = xelatex
% !TEX encoding = UTF-8 Unicode
% !Mode:: "TeX:UTF-8"

\documentclass{resume}
\usepackage{zh_CN-Adobefonts_external} % Simplified Chinese Support using external fonts (./fonts/zh_CN-Adobe/)
% \usepackage{NotoSansSC_external}
% \usepackage{NotoSerifCJKsc_external}
% \usepackage{zh_CN-Adobefonts_internal} % Simplified Chinese Support using system fonts
\usepackage{linespacing_fix} % disable extra space before next section
\usepackage{cite}

\begin{document}
\pagenumbering{gobble} % suppress displaying page number

\name{訾宇}

\basicInfo{
  \email{zyu539@gmail.com} \textperiodcentered\
  \phone{(+1) 412-482-0674} \textperiodcentered\
  \linkedin[billryan8]{https://www.linkedin.com/in/yu-zi-15a001123/}}
 
\section{\faGraduationCap\  教育背景}
\datedsubsection{\textbf{卡内基梅隆大学, 信息系统管理}, 美国}{2021 -- 2022}
\textit{硕士, 优秀毕业生(Distinction of Heinz College)}
\datedsubsection{\textbf{奥克兰大学, 软件工程}, 新西兰}{2014 -- 2018}
\textit{学士, 工程学院院长嘉许名单(Dean's Honours List), First Class}

\section{\faUsers\ 工作经历}
\datedsubsection{\textbf{Google LLC (谷歌)} 得克萨斯}{2023年5月 -- 至今}
\role{软件工程师(Java)}
任职于 Google Cloud 云支持平台的 CaseCentral 团队。负责构建相关系统，使客户和合作伙伴能够联系 Google 获取支持，旨在跨渠道和产品线提供一致且无缝的用户支持体验。
\begin{itemize}
  \item 构建并维护一套 API，用于管理客户支持相关的实体对象，如案例 (Case)、评论 (Comment) 和附件 (Attachment) 等
  \item 设计并构建客户支持相关实体的Database Schema，将相关数据从Salesforce迁移至Google内部数据库并保持同步
  \item 持续整合来自其他 Google Cloud 和 Workspace 内部产品线的支持平台
\end{itemize}

\datedsubsection{\textbf{Amazon.com, Inc. (亚马逊)}}{2022年10月 -- 2023年5月}
\role{软件工程师(Java)}
\begin{onehalfspacing}
任职于购物体验部门下的BrandFollow团队。主要职责为开发并维护前后端服务，使客户能够关注品牌并获得独家权益，包括品牌推荐和独家优惠
\begin{itemize}
  \item 开发并维护用于管理亚马逊全平台品牌标识符和元数据的服务，作为品牌作为可关注实体的“唯一事实来源” (Source of Truth)
  \item 构建服务间的数据迁移管道，同时确保数据质量和服务可用性
\end{itemize}
\end{onehalfspacing}

\datedsubsection{\textbf{Kakapo Technologies}}{2018年4月 -- 2021年5月}
\role{软件工程师(Python)}
\begin{onehalfspacing}
为澳大利亚国民银行 (National Australia Bank) 开发并维护服务，用于协调不同系统间的交易信息
\begin{itemize}
  \item 开发并维护对账系统定期从各系统提取数据，识别本应一致的交易，并定位其中的差异
  \item 开发运行于AWS的服务，监控来自不同系统的数据到达情况，并通过消息队列协调包含数千个对账工作流
  \item 开发解析器和构建器，处理和分析数百种不同格式的数据
\end{itemize}
\end{onehalfspacing}

% Reference Test
\datedsubsection{\textbf{Cloud-based Anomaly Route Detection System}}{2017}
Build a mobile application which can collect real-time trajectory data to analyze if taxi drivers are taking an anomalous route.
\begin{itemize}
 \item A cloud-based Java backend to store and normalize route data, train a prediction model periodically, \linebreak and justify the current route with the trained model, achieving a correctness rate of over 90\%
 \item ETL process to clean up \& feature engineering trajectory data
\end{itemize}

\section{\faCogs\ 技能}
% increase linespacing [parsep=0.5ex]
\begin{itemize}[parsep=0.5ex]
  \item 编程语言: Java, Python, Go, C, TypeScript
  \item 技能: 微服务(Micro Services), gRPC, 数据库(DynanmoDB, Spanner, MySQL等), 机器学习， Kubernetes
  \item 兴趣方向: 云计算, 分布式系统, 大语言模型
  \item 语言: 英语 - 熟练(TOEFL 107， GRE 332)
\end{itemize}

\section{\faHeartO\ 获奖情况}
\datedline{\textit{第11名}, ACM-ICPC 国际大学生程序设计竞赛 南太平洋赛区决赛}{2016}
\datedline{\textit{第2名}, 新西兰程序设计竞赛 (New Zealand Programming Contest)}{2016}
\datedline{\textit{第87名}, IEEEXtreme 24小时极限编程竞赛}{2016}
\datedline{\textit{第3名}, 新西兰程序设计竞赛 (New Zealand Programming Contest)}{2015}

%% Reference
%\newpage
%\bibliographystyle{IEEETran}
%\bibliography{mycite}
\end{document}
